\documentclass[11pt]{article}

\usepackage[margin=1in]{geometry}
\usepackage{array}
\usepackage{amsmath}
\usepackage{amssymb}
\usepackage{booktabs}
\usepackage{enumitem}
\usepackage{xcolor}
\usepackage{hyperref}
\usepackage{longtable}
\usepackage{listings}

\hypersetup{
  colorlinks=true,
  linkcolor=blue!60!black,
  urlcolor=blue!60!black
}

\lstdefinelanguage{Lean}{
  keywords={theorem,lemma,def,abbrev,structure,where,by,have,show,exact,
    refine,obtain,let,fun,if,then,else,match,with,forall,exists,Prop,Type},
  sensitive=true,
  comment=[l]{--},
  morecomment=[s]{/-}{-/},
  morestring=[b]"
}

\lstset{
  language=Lean,
  basicstyle=\ttfamily\small,
  columns=fullflexible,
  breaklines=true,
  keepspaces=true,
  frame=single,
  framerule=0.3pt,
  backgroundcolor=\color{black!3}
}

\newcommand{\LeanTODO}[1]{%
  \begin{quote}
  \ttfamily\small\raggedright #1
  \end{quote}
}

\newcommand{\BenchHead}[1]{%
  \medskip\noindent\textbf{#1}\par\noindent
}

\newcommand{\BenchSubhead}[1]{%
  \smallskip\noindent\textit{#1}\par\noindent
}

\newcommand{\Exercise}[7]{%
  \subsection{#1}
  \begin{description}[style=nextline,leftmargin=3.8cm,font=\normalfont\bfseries]
    \item[Tier] #2
    \item[Theme] #3
    \item[Exercise prompt] #4
    \item[Benchmark intent] #5
    \item[Expected dependencies] #6
    \item[Notes] #7
  \end{description}

  \BenchHead{Layer 1: natural-language discovery run}
  This run measures whether the model can turn the exercise prompt into a useful
  Lean theorem and proof.

  \BenchSubhead{Condition A setup}
  Exercise prompt plus only the neutral Lean scaffold needed to state the task:
  basic types, notation, and algorithm definitions if required.  No
  ComputationalMathematics imports, source ledger, stability bounds, proof lemmas, or
  proof logic.

  \BenchSubhead{Condition A generated theorem and proof}
  \LeanTODO{-- TODO: paste model output from natural-language Condition A run}

  \BenchSubhead{Condition C setup}
  Same exercise prompt, but with ComputationalMathematics imports, source-oriented library
  documentation, and available project context.

  \BenchSubhead{Condition C generated theorem and proof}
  \LeanTODO{-- TODO: paste model output from natural-language Condition C run}

  \BenchSubhead{Discovery assessment}
  TODO: record whether the two generated Lean statements prove the same
  mathematical theorem, whether either statement is too weak or has hidden
  assumptions, and whether the result is eligible for the fixed-theorem run.

  \BenchHead{Layer 2: fixed Lean theorem run}
  This run measures proof assistance only.  The canonical theorem below must be
  frozen before either condition is run.

  \BenchSubhead{Canonical Lean theorem statement}
  \LeanTODO{-- TODO: theorem or lemma statement used verbatim in both fixed-theorem conditions}

  \BenchSubhead{Minimal statement context}
  \LeanTODO{-- TODO: neutral stubs/imports needed to state the theorem; no bounds or proof lemmas}

  \BenchSubhead{Condition A proof}
  \LeanTODO{-- TODO: generated proof from fixed Lean theorem with minimal statement scaffold}

  \BenchSubhead{Condition C library context}
  \LeanTODO{-- TODO: ComputationalMathematics imports/docs/lemmas made visible for Condition C}

  \BenchSubhead{Condition C proof}
  \LeanTODO{-- TODO: generated proof from the exact same fixed Lean theorem with library context}

  \BenchSubhead{Fixed-theorem assessment}
  TODO: record success or failure, proof length, wall time, use of source-matched
  library lemmas, hidden assumptions, remaining \texttt{sorry}s, and whether the
  proof is implementation-backed or contract-transfer.

  \BenchSubhead{Natural-language explanation of accepted proof}
  TODO.
}

\title{Higham Exercise Benchmark Candidates}
\author{Lean Computational Mathematics Benchmarking}
\date{\today}

\begin{document}
\maketitle

\section{Purpose}

This document collects candidate exercises for the Lean Computational Mathematics benchmark.
It is a benchmark ledger, not a source-coverage ledger.

The benchmark goal is to test whether access to the Lean Computational Mathematics library
makes formalizing floating-point stability-analysis tasks faster and more
reliable than working from a minimal statement scaffold.  Preferred exercises ask for
perturbation bounds, backward-error bounds, forward-error bounds, residual
bounds, condition numbers, or algorithmic stability statements.  We avoid
tasks that are only lookup, empirical exploration, or exact algebra unless
they directly support a stability bound.

Long verbatim exercise statements are intentionally not copied here.  Each
entry records the exercise number, a source-faithful prompt paraphrase with
the essential equations or bound target, and placeholders for the Lean task
statement and generated proofs.

The early-chapter numbering follows the user-curated exercise list.  In a few
places the local PDF extraction differs slightly from the screenshot numbering;
the prompt text identifies the intended exercise content.

\section{Benchmark Design}

\begin{description}[style=nextline,leftmargin=3.2cm,font=\normalfont\bfseries]
  \item[Condition A] Minimal statement scaffold.  This includes only what is
  needed to state the exercise in Lean in the same form as Condition C: basic
  types, notation, local structures, and possibly concrete algorithm definitions.
  It must not include stability bounds, source-derived proof lemmas, or proof
  logic.
  \item[Condition C] The same task file and the same neutral statement scaffold,
  plus the Lean Computational Mathematics library context and documentation.
\end{description}

Each exercise is evaluated in two layers.

\paragraph{Layer 1: natural-language discovery}
Both conditions receive the exercise prompt and, if needed, the same neutral
statement scaffold.  Condition A has no stability-library help, while Condition
C sees the Lean Computational Mathematics library and documentation.  This tests whether library
context helps the model choose a good Lean formalization and prove it.  The
result is inherently noisy because the two runs may choose different theorem
statements.

\paragraph{Layer 2: fixed Lean theorem}
A canonical Lean theorem statement is prepared and frozen before benchmarking.
Both conditions receive exactly this theorem.  Condition A receives the minimal
statement scaffold only; Condition C receives the same theorem and scaffold plus
Lean Computational Mathematics context.  This is the primary benchmark because it tests proof
support without theorem-selection drift.

For both layers, paste the generated proof under the corresponding condition.
For Layer 1, also record whether the generated Lean statements are genuinely
the same mathematical claim.  For Layer 2, the theorem statement must be byte-for-byte
identical across conditions.

\section{Code Block Convention}

Use Lean listings for task statements and proofs:

\begin{lstlisting}
theorem example_task : True := by
  trivial
\end{lstlisting}

\section{Current Library Coverage Notes}

\begin{itemize}
  \item Rounded \texttt{exp}, rounded \texttt{log}, and \texttt{log1p} models are
  not yet a settled part of the library.  Exercises that need these operations
  should remain in the benchmark: they can show whether the model can introduce
  adequate local stubs, whether a proof exposes a real library extension need,
  or whether the current library context is insufficient.  If a run scopes such
  a task to exact algebra, that scoping must be recorded explicitly in the
  assessment rather than treated as the original task.
\end{itemize}

\section{Candidate Index}

\small
\begin{longtable}{>{\raggedright\arraybackslash}p{0.07\textwidth}
                  >{\raggedright\arraybackslash}p{0.09\textwidth}
                  >{\raggedright\arraybackslash}p{0.08\textwidth}
                  >{\raggedright\arraybackslash}p{0.19\textwidth}
                  >{\raggedright\arraybackslash}p{0.37\textwidth}}
\toprule
Chapter & Exercise & Tier & Theme & Benchmark intent \\
\midrule
\endfirsthead
\toprule
Chapter & Exercise & Tier & Theme & Benchmark intent \\
\midrule
\endhead
1 & 1.9 & Best & Forward and backward error & Cramer's rule for 2 by 2 systems; derive forward and residual bounds. \\
1 & 1.10 & Good & Forward error & Two-pass sample variance error bound. \\
1 & 1.7 & Good & Condition number & Componentwise and normwise condition numbers for sample variance. \\
1 & 1.5 & Maybe & Accurate reformulation & Large-\(n\) computation of \((1+1/n)^n\); useful probe for rounded log/exp support. \\
3 & 3.5 & Best & Backward error & Matrix multiplication backward error with perturbation in either factor. \\
3 & 3.2 & Good & Gamma bound & Sharpened product-error gamma bound under special hypotheses. \\
3 & 3.8 & Good & Matrix product forward error & Bound for product of a sequence of matrices. \\
3 & 3.12 & Good & Quadrature error & Floating-point summation and function evaluation error bound for quadrature. \\
3 & 3.7 & Maybe & Complex backward error & Complex analogues of real backward-error results. \\
4 & 4.3 & Best & Summation forward error & Refined recursive-summation error representation and ordering bound. \\
4 & 4.7 & Maybe & Stability comparison & Analyze summation via \(\log(\prod_i \exp x_i)\); useful probe for rounded exp/log support. \\
4 & 4.8 & Maybe & Grid generation error & Compare recurrence and direct formulas for equally spaced grid points. \\
5 & 5.2 & Best & Polynomial evaluation error & Error analysis of beginner polynomial evaluation. \\
5 & 5.3 & Best & Horner-style error bound & Error bound for even/odd splitting with Horner on each part. \\
5 & 5.6 & Best & Polynomial backward/forward error & Bound for computed polynomial from Horner-like evaluation. \\
5 & 5.5 & Good & Product-form evaluation & Error analysis for root-product evaluation of a polynomial. \\
7 & 7.2 & Good & Residual-to-forward error & Relate residual, forward error, and condition number for linear systems. \\
7 & 7.7 & Good & Backward error comparison & Compare backward-error quantities under perturbed right-hand sides. \\
8 & 8.3 & Best & Triangular solve forward error & Bound components of triangular solve under unit upper triangular assumptions. \\
8 & 8.4 & Best & Componentwise conditioning & Bound triangular M-matrix solve condition measure. \\
8 & 8.7 & Maybe & Diagonal dominance bound & Use diagonal dominance to bound inverse norms; use scoped part only. \\
9 & 9.6 & Good & Growth factor & Totally nonnegative matrices and growth factor one. \\
9 & 9.8 & Good & LU backward structure & Inverse totally nonnegative case and absolute LU factorization. \\
9 & 9.9 & Good & Growth factor bound & Bound GE growth factor without pivoting using \(\lvert L\rvert\lvert U\rvert\). \\
9 & 9.13 & Good & Threshold pivoting & Growth-factor bound for sparse threshold pivoting. \\
9 & 9.14 & Maybe & Pivoting equivalence & Pre-pivoted GEPP equivalences; mostly structural. \\
19 & 19.2 & Best & Orthogonality error & Bound \(\widehat P^T\widehat P-I\) for a computed Householder reflector. \\
19 & 19.6 & Best & Householder update error & Derive one-stage Householder QR update error and row-scaling interpretation. \\
19 & 19.10 & Good & Gram-Schmidt instability & Analyze orthogonality of CGS and MGS on a small Lauchli-type matrix. \\
19 & 19.12 & Good & Backward-error structure & Convert block perturbation data into a QR backward-error statement. \\
20 & 20.2 & Best & Least-squares backward stability & Prove or scope-formalize Householder QR LS backward-stability theorem. \\
20 & 20.5 & Best & MGS LS backward stability & Analogue of the QR LS backward-stability result for MGS. \\
20 & 20.6 & Best & Residual error & Bound difference between true and computed LS residual norms for normal equations. \\
20 & 20.8 & Good & LS backward error & Special case of LS backward-error formula when the right-hand side is unperturbed. \\
20 & 20.9 & Good & Stable formula & Derive cancellation-resistant formula for normwise LS backward error. \\
20 & 20.11 & Good & First-order approximation & Show LSE condition expression reduces to unconstrained LS first-order form. \\
22 & 22.8 & Best & Bidiagonal inverse perturbation & Derive exact inverse perturbation formula and componentwise bound. \\
22 & 22.11 & Best & Structured condition number & Derive structured condition number for primal Vandermonde systems. \\
22 & 22.3 & Good & Algorithm stability & Scoped residual or forward-error task for Vandermonde inverse algorithm. \\
22 & 22.7 & Good & Conditioning & Condition bounds for Chebyshev-Vandermonde matrices at special nodes. \\
23 & 23.6 & Best & Fast matrix multiply error & Prove 3M plus Strassen normwise error bound. \\
24 & 24.1 & Best & FFT convolution error & Bound DFT convolution rounding error and integer recovery condition. \\
25 & 25.1 & Good & Perturbed iteration stability & Error floor for contraction iteration with bounded perturbations. \\
27 & 27.5 & Good & Robust norm computation & Correctness and stability properties of scaled 2-norm algorithm. \\
27 & 27.6 & Good & Robust Pythagorean sum & Halley/pythag iteration, convergence, and robustness properties. \\
\bottomrule
\end{longtable}
\normalsize

\section{Exercise Workspaces}

\Exercise{Ch01 Ex. 1.9 -- Cramer's Rule Stability}
{Best}
{Forward and residual error}
{For a \(2\) by \(2\) system \(Ax=b\), Cramer's rule forms
\[
d=a_{11}a_{22}-a_{21}a_{12},\quad
x_1=(b_1a_{22}-b_2a_{12})/d,\quad
x_2=(a_{11}b_2-a_{21}b_1)/d .
\]
Assuming \(d\) is exact, show that the computed \(\widehat x\) satisfies
\[
\frac{\lVert x-\widehat x\rVert_\infty}{\lVert x\rVert_\infty}
\le \gamma_3\,\operatorname{cond}(A,x),\qquad
\lVert b-A\widehat x\rVert_\infty
\le \gamma_3\,\operatorname{cond}(A^{-1})\lVert b\rVert_\infty ,
\]
with the componentwise condition quantities defined using
\(\lvert A^{-1}\rvert\lvert A\rvert\lvert x\rvert\).}
{Derive forward and residual bounds from the stated operation model.}
{Floating-point product/sum error; small linear-system perturbation bounds.}
{Good first benchmark because the task is small but not a lookup.}

\Exercise{Ch01 Ex. 1.10 -- Two-Pass Sample Variance}
{Good}
{Forward error}
{For the two-pass sample variance formula \(V(x)\), prove that the computed
quantity \(\widehat V=\operatorname{fl}(V(x))\) obeys
\[
\frac{\lvert V-\widehat V\rvert}{V}\le (n+3)u+O(u^2).
\]
The point is that the leading term should not depend on the sample-variance
condition numbers from Ex. 1.7.}
{Show the relative error is bounded by a linear term in \(n u\) plus higher-order terms.}
{Summation error; dot/product error; variance definitions.}
{More valuable after sample-variance and summation APIs are settled.}

\Exercise{Ch01 Ex. 1.7 -- Sample Variance Condition Numbers}
{Good}
{Condition number}
{For sample variance \(V(x)\), define a componentwise condition number using
\(\lvert \Delta x_i\rvert\le \epsilon\lvert x_i\rvert\) and show
\[
\kappa_C
=\frac{2\sum_i\lvert x_i-\bar x\rvert\lvert x_i\rvert}
{(n-1)V(x)} .
\]
For the normwise perturbation set \(\lVert\Delta x\rVert_2\le
\epsilon\lVert x\rVert_2\), show
\[
\kappa_N=\frac{2\lVert x\rVert_2}{\sqrt{(n-1)V(x)}}
=2\left(1+\frac{n}{n-1}\frac{\bar x^2}{V(x)}\right)^{1/2}
\ge \kappa_C .
\]}
{Formalize the condition-number formulas and comparison.}
{Exact calculus or first-order perturbation framework; variance definitions.}
{Useful companion to Ch01 Ex. 1.10.}

\Exercise{Ch01 Ex. 1.5 -- Accurate Large-n Power}
{Maybe}
{Accurate reformulation}
{First obtain an accurate method for \(\log(1+x)\) for all \(x>-1\), including
small \(\lvert x\rvert\), assuming the library log has relative error at most
\(u\).  Then use
\[
(1+1/n)^n=\exp(n\log(1+1/n))
\]
to compute \((1+1/n)^n\) accurately for large \(n\).}
{Test whether the run introduces adequate rounded log/exp stubs or exposes a
library extension need.}
{Rounded log/exp models; possible local \texttt{log1p} specification.}
{Keep in the benchmark as a coverage probe.  If a run proves only an exact
algebra variant, record that explicitly as a scoped result.}

\Exercise{Ch03 Ex. 3.5 -- Matrix Multiplication Backward Error}
{Best}
{Backward error}
{For nonsingular \(A,B\in\mathbb R^{n\times n}\), show that the computed product
can be written as
\[
\operatorname{fl}(AB)=(A+\Delta A)B,\qquad
\lvert\Delta A\rvert\le \gamma_n\lvert A\rvert\lvert B\rvert\lvert B^{-1}\rvert .
\]
Also derive the corresponding backward-error representation in which \(B\) is
perturbed instead of \(A\).}
{Prove a componentwise or normwise backward-error statement for matrix multiplication.}
{Dot-product backward error; matrix-vector or matrix-matrix composition.}
{Excellent core benchmark.}

\Exercise{Ch03 Ex. 3.2 -- Sharpened Product Gamma Bound}
{Good}
{Gamma bound}
{In Lemma 3.1, specialize to the case \(\rho_i\equiv 1\).  Prove the sharper
product-error estimate
\[
\lvert\theta_n\rvert\le \frac{nu}{1-\frac12nu},\qquad nu<2 .
\]}
{Prove the strengthened gamma inequality.}
{Gamma algebra and product-of-errors lemmas.}
{Calibration task for gamma infrastructure.}

\Exercise{Ch03 Ex. 3.8 -- Matrix Sequence Product Error}
{Good}
{Matrix product forward error}
{For \(A_1,\ldots,A_k\in\mathbb R^{n\times n}\), bound the forward error in
the computed matrix product by proving a statement of the form
\[
\lVert A_1\cdots A_k-\operatorname{fl}(A_1\cdots A_k)\rVert_F
\le (kn^2u+O(u^2))\lVert A_1\rVert_2\cdots\lVert A_k\rVert_2 .
\]}
{Show a Frobenius or normwise error bound for a finite matrix product.}
{Matrix multiplication error and norm composition.}
{Good if scoped to fixed product length or a clean recursive product definition.}

\Exercise{Ch03 Ex. 3.12 -- Quadrature Rule Error}
{Good}
{Quadrature error}
{For a quadrature rule
\[
I(f)=\int_a^b f(x)\,dx \approx J(f)=\sum_{i=1}^n w_i f(x_i),
\]
assume \(w_i,x_i\) are floating-point data, the weighted sum is evaluated
left-to-right, and
\(\operatorname{fl}(f(x_i))=f(x_i)(1+\eta_i)\), \(\lvert\eta_i\rvert\le\eta\).
Derive and interpret a bound for
\(\lvert I(f)-\widehat J(f)\rvert\), where
\(\widehat J(f)=\operatorname{fl}(J(f))\).}
{Derive and interpret a bound for \(\lvert I(f)-\widehat J(f)\rvert\).}
{Summation bounds; weighted sum bounds.}
{Good applied stability task.}

\Exercise{Ch03 Ex. 3.7 -- Complex Backward Error Analogues}
{Maybe}
{Complex backward error}
{Give complex analogues of the real backward-error results for inner products
and matrix-vector products, replacing transposes by conjugate transposes where
appropriate and using complex modulus in the componentwise bounds.}
{Prove one specific complex dot-product or matrix-vector backward-error analogue.}
{Complex absolute values; real/complex transfer lemmas.}
{Avoid using it as a broad prose task.}

\Exercise{Ch04 Ex. 4.3 -- Recursive Summation Refined Bound}
{Best}
{Summation forward error}
{Let \(S_n=\sum_{i=1}^n x_i\) be computed by natural-order recursive summation.
Show the refined representation
\[
\widehat S_n=(x_1+x_2)(1+\theta_{n-1})
+\sum_{i=3}^n x_i(1+\theta_{n-i+1}),\qquad
\lvert\theta_k\rvert\le\gamma_k=\frac{ku}{1-ku}.
\]
Deduce
\[
\lvert \widehat S_n-S_n\rvert
\le (\lvert x_1\rvert+\lvert x_2\rvert)\gamma_{n-1}
+\sum_{i=3}^n \lvert x_i\rvert\gamma_{n-i+1},
\]
and determine the ordering of the \(x_i\) that minimizes this bound.}
{Prove the displayed weighted gamma bound and ordering consequence.}
{Recursive summation backward error and gamma monotonicity.}
{Strong because the proof is stability-specific and not just algebra.}

\Exercise{Ch04 Ex. 4.7 -- Log-Exp Summation}
{Maybe}
{Stability comparison}
{Analyze the accuracy of computing a sum through the identity
\[
S_n=\log\left(\prod_{i=1}^n e^{x_i}\right).
\]
The intended task is to compare this with direct summation and expose the
rounding-error risks introduced by rounded \(\exp\), product accumulation, and
rounded \(\log\).}
{Test whether the run can model rounded exp/log plus product accumulation, or
identify the missing operation model cleanly.}
{Rounded exp/log models; product error.}
{Keep in the benchmark as a coverage probe.  If a run proves only an exact
algebra variant, record that explicitly as a scoped result.}

\Exercise{Ch04 Ex. 4.8 -- Grid Point Formation}
{Maybe}
{Grid generation error}
{For equally spaced points on \([a,b]\),
\[
x_i=a+ih,\qquad h=(b-a)/n,\qquad i=0,\ldots,n,
\]
compare the accuracy of forming \(x_i\) by repeated addition
\(x_i=x_{i-1}+h\), by the direct formula \(a+ih\), and by the convex
combination \(a(1-i/n)+(i/n)b\).  Assume \(a,b\) are floating-point numbers
but \(h\) need not be.}
{Prove one scoped error comparison.}
{Rounded addition/multiplication and simple recurrences.}
{Keep the task narrow.}

\Exercise{Ch05 Ex. 5.2 -- Beginner Polynomial Evaluation}
{Best}
{Polynomial evaluation error}
{For \(p(x)=a_0+a_1x+\cdots+a_nx^n\), analyze the direct evaluator
\[
q:=a_0,\quad y:=1,\quad
\text{for }i=1:n\text{ do } y:=xy,\ q:=q+a_i y .
\]
Derive a rounding-error bound for the returned \(q\).}
{Prove a coefficientwise or normwise error bound.}
{Product recurrence; summation error; polynomial definitions.}
{Good contrast with Horner tasks.}

\Exercise{Ch05 Ex. 5.3 -- Even/Odd Polynomial Splitting}
{Best}
{Horner-style error bound}
{Let \(p(x)=a_0+a_1x+\cdots+a_nx^n\) with \(n=2m\).  Write
\[
p(x)=a_0+a_2y+\cdots+a_{2m}y^m
+x(a_1+a_3y+\cdots+a_{2m-1}y^{m-1}),\qquad y=x^2 .
\]
Obtain an error bound for \(\operatorname{fl}(p(x))\) when each part is evaluated
by Horner's rule and the final combination is rounded.}
{Derive an error bound for the split evaluation.}
{Horner backward error; product error; polynomial splitting.}
{Strong because it composes existing stability lemmas.}

\Exercise{Ch05 Ex. 5.6 -- Computed Polynomial Bound}
{Best}
{Polynomial backward/forward error}
{For the polynomial \(P_3\) in the surrounding chapter, prove that the computed
polynomial \(\widehat P_3\) produced by Horner's method satisfies
\[
\lVert P_3-\widehat P_3\rVert
\le n(m+1)u\,\widetilde p_3(\lVert X\rVert)+O(u^2),
\qquad
\widetilde p_3(x)=\sum_k \lVert A_k\rVert x^k ,
\]
for the \(1\)-norm or \(\infty\)-norm.}
{Formalize the norm bound involving a nonnegative majorant polynomial.}
{Horner coefficientwise backward error and norm bounds.}
{Strong higher-level Horner benchmark.}

\Exercise{Ch05 Ex. 5.5 -- Root Product Evaluation}
{Good}
{Product-form evaluation}
{If \(p(x)=\sum_{i=0}^n a_ix^i\) has roots \(x_1,\ldots,x_n\), evaluate it by
the product form
\[
p(x)=a_n\prod_{i=1}^n(x-x_i).
\]
Give a rounding-error analysis for this product evaluation.}
{Derive a product-form backward or forward error bound.}
{Product rounding lemmas.}
{Good if roots and polynomial product form are already available.}

\Exercise{Ch07 Ex. 7.2 -- Residual and Forward Error}
{Good}
{Residual-to-forward error}
{Let \(Ax=b\), \(A\in\mathbb R^{n\times n}\), and let \(y\) be any approximate
solution with residual \(r=b-Ay\).  For any subordinate matrix norm, prove
\[
\frac{\lVert r\rVert}{\lVert A\rVert\lVert x\rVert}
\le
\frac{\lVert x-y\rVert}{\lVert x\rVert}
\le
\kappa(A)\frac{\lVert r\rVert}{\lVert A\rVert\lVert x\rVert}.
\]
Interpret the bound.}
{Prove the two-sided residual/forward-error inequality.}
{Normed matrix-vector inequalities and inverse bounds.}
{Good foundational perturbation task.}

\Exercise{Ch07 Ex. 7.7 -- Backward Error Comparison}
{Good}
{Backward error comparison}
{For the normwise and componentwise backward-error measures from the chapter,
prove the comparison inequalities
\[
\omega_{\lvert A\rvert,\lvert b\rvert}(y)
\le \omega_{\lvert A\rvert,0}(y)
\le \frac{2\omega_{\lvert A\rvert,\lvert b\rvert}(y)}
{1-\omega_{\lvert A\rvert,\lvert b\rvert}(y)}
\]
and
\[
\eta_{A,b}(y)\le \eta_{A,0}(y)
\le \frac{2\eta_{A,b}(y)}{1-\eta_{A,b}(y)} .
\]}
{Derive the displayed comparison inequalities.}
{Backward-error definitions and algebraic manipulation of bounds.}
{Useful for library-level error-measure definitions.}

\Exercise{Ch08 Ex. 8.3 -- Triangular Solve Component Bound}
{Best}
{Triangular solve forward error}
{Let \(U\in\mathbb R^{n\times n}\) be unit upper triangular with
\(\lvert u_{ij}\rvert\le 1\) for \(j>i\).  Using the triangular-solve theorem,
show that the computed substitution solution \(\widehat x\) to \(Ux=b\) satisfies
\[
\lvert x_i-\widehat x_i\rvert
\le 2^{\,n-i}\bigl((n^2+n+1)u+O(u^2)\bigr)\lVert b\rVert_\infty .
\]
Compare this with the result obtained from the more general triangular-solve
bound.}
{Prove the componentwise bound for unit upper triangular systems with bounded off-diagonals.}
{Triangular solve backward/forward bounds.}
{Strong direct stability task.}

\Exercise{Ch08 Ex. 8.4 -- Triangular M-Matrix Conditioning}
{Best}
{Componentwise conditioning}
{Let \(T\in\mathbb R^{n\times n}\) be triangular, assume \(T=M(T)\), and suppose
\(Tx=b\ge 0\).  Prove
\[
\lvert T^{-1}\rvert\lvert T\rvert\lvert x\rvert\le (2n-1)\lvert x\rvert,
\]
and hence \(\operatorname{cond}(T,x)\le 2n-1\).}
{Formalize the bound \(\operatorname{cond}(T,x)\le 2n-1\) under the stated hypotheses.}
{Triangular comparison and M-matrix inverse bounds.}
{Strong if the M-matrix definitions are stable.}

\Exercise{Ch08 Ex. 8.7 -- Diagonal Dominance Inverse Bound}
{Maybe}
{Diagonal dominance bound}
{For a strictly row diagonally dominant matrix satisfying
\[
\alpha_i=\lvert a_{ii}\rvert-\sum_{j\ne i}\lvert a_{ij}\rvert>0,
\]
prove
\[
\lVert A^{-1}\rVert_\infty\le \frac{1}{\min_i\alpha_i}.
\]
Optionally extend this to the scaled case \(AD\) with positive diagonal
\(D=\operatorname{diag}(d_i)\), giving
\(\lVert A^{-1}\rVert_\infty\le \lVert D\rVert_\infty/\min_i\beta_i\).}
{Use only a scoped subpart, preferably the inverse-norm bound.}
{Matrix norm and diagonal dominance lemmas.}
{More exact matrix analysis than floating-point stability.}

\Exercise{Ch09 Ex. 9.6 -- Totally Nonnegative Growth Factor}
{Good}
{Growth factor}
{If \(A\in\mathbb R^{n\times n}\) is nonsingular and totally nonnegative, show
that it has an LU factorization \(A=LU\) with \(L\ge 0\) and \(U\ge 0\).  Using
the determinant inequality for totally nonnegative matrices, deduce that Gaussian
elimination without pivoting has growth factor \(\rho_n=1\).}
{Formalize the growth-factor consequence.}
{LU factorization; total nonnegativity; determinant inequalities.}
{Specialized but relevant to GE stability.}

\Exercise{Ch09 Ex. 9.8 -- Absolute LU Factorization}
{Good}
{LU backward structure}
{If \(A\in\mathbb R^{n\times n}\) is nonsingular and \(A^{-1}\) is totally
nonnegative, show that \(A\) has an LU factorization \(A=LU\) satisfying
\[
\lvert A\rvert=\lvert L\rvert\lvert U\rvert .
\]}
{Derive the LU absolute-value relation and stability implication.}
{LU exact algebra and sign-structure lemmas.}
{Good support for growth-factor and stability reasoning.}

\Exercise{Ch09 Ex. 9.9 -- GE Growth Bound Without Pivoting}
{Good}
{Growth factor bound}
{For Gaussian elimination without pivoting, prove the growth-factor estimate
\[
\rho_n\le 1+\frac{n\lVert \lvert L\rvert\lvert U\rvert\rVert_\infty}
{\lVert A\rVert_\infty}.
\]}
{Prove the growth-factor inequality.}
{GE/LU definitions and max-norm inequalities.}
{Good compact task if growth factor is already defined.}

\Exercise{Ch09 Ex. 9.13 -- Threshold Pivoting Growth Bound}
{Good}
{Threshold pivoting}
{In sparse Gaussian elimination with threshold pivoting, a pivot in column \(k\)
is chosen so that
\[
\lvert a^{(k)}_{ik}\rvert\ge \tau\max_{m\ge k}\lvert a^{(k)}_{mk}\rvert,\qquad
\tau\in(0,1].
\]
Show that
\[
\max_i\lvert a^{(k)}_{ij}\rvert
\le (1+\tau^{-1})^{\mu_j}\max_i\lvert a_{ij}\rvert,
\]
where \(\mu_j\) is the number of nonzeros in the \(j\)th column of \(U\), and
deduce a growth-factor bound.}
{Prove the per-column growth recurrence and resulting growth-factor bound.}
{Sparse column counts; threshold pivoting definitions; growth-factor algebra.}
{Good later benchmark; definitions may need local setup.}

\Exercise{Ch09 Ex. 9.14 -- Pre-Pivoted GEPP Equivalence}
{Maybe}
{Pivoting equivalence}
{Let \(A\) be pre-pivoted for GEPP, meaning GEPP performs no row interchanges
on \(A\).  Show GEPP computes the same LU factorization as the first pivoting
method from the chapter and as pairwise pivoting applied to \(\Pi A\), where
\(\Pi\) is the row-reversal permutation, using the natural strategy that keeps
multipliers bounded by \(1\).}
{Use only if scoped to a stability consequence.}
{Permutation matrices and GEPP definitions.}
{Mostly structural, so lower priority.}

\Exercise{Ch19 Ex. 19.2 -- Computed Householder Orthogonality}
{Best}
{Orthogonality error}
{Let \(\widehat P=I-\widehat\beta\,\widehat v\,\widehat v^T\), where
\(\widehat\beta\) and \(\widehat v\) are the computed Householder quantities from
Lemma 19.1.  Derive a bound for
\[
\lVert \widehat P^T\widehat P-I\rVert_2 .
\]}
{Prove a bound for \(\widehat P^T\widehat P-I\) using the computed beta/vector error model.}
{Householder reflector definitions and norm inequalities.}
{Very good QR benchmark.}

\Exercise{Ch19 Ex. 19.6 -- Householder QR One-Step Error}
{Best}
{Householder update error}
{In Householder QR, prove the vector-size bound for the \(k\)th Householder
vector,
\[
\sqrt2\,\lVert a_k^{(k)}(k:m)\rVert_2
\le \lVert v_k\rVert_2
\le 2\lVert a_k^{(k)}(k:m)\rVert_2 .
\]
Then, for \(j>k\), analyze the computed update
\(\widehat a_j^{(k+1)}=\operatorname{fl}(\widehat P_k\widehat a_j^{(k)})\),
where \(\widehat P_k=I-\beta_k\widehat v_k\widehat v_k^T\) and
\(\lvert\Delta v_k\rvert\le\widetilde\gamma_{m-k}\lvert v_k\rvert\).  Show
\[
\widehat a_j^{(k+1)}=P_k\widehat a_j^{(k)}+f_j^{(k)},
\]
with the leading zero pattern in \(f_j^{(k)}\) and the displayed componentwise
bound involving \(u\lvert\widehat a_j^{(k)}\rvert\) and
\(\widetilde\gamma_{m-k}\lvert v_k\rvert\).}
{Prove the residual form and componentwise/normwise bound for a panel update.}
{Householder vector bounds; reflector application error.}
{Excellent because it composes several QR facts.}

\Exercise{Ch19 Ex. 19.10 -- CGS/MGS Orthogonality Example}
{Good}
{Gram-Schmidt instability}
{For the Lauchli-type matrix
\[
A=\begin{bmatrix}
1&1&1\\
\epsilon&0&0\\
0&\epsilon&0\\
0&0&\epsilon
\end{bmatrix},
\]
assuming \(\operatorname{fl}(1+\epsilon^2)=1\), evaluate the \(Q\) matrices
produced by classical Gram-Schmidt and modified Gram-Schmidt, and assess their
orthonormality.}
{Formalize the computed \(Q\) matrices under a rounding assumption and bound orthogonality loss.}
{Gram-Schmidt algorithm definitions.}
{Good as an example task, but less general than 19.6.}

\Exercise{Ch19 Ex. 19.12 -- QR Backward-Error Structure}
{Good}
{Backward-error structure}
{Assume matrices satisfy
\[
\begin{bmatrix}A+\Delta A_1\\ \Delta A_2\end{bmatrix}
=
\begin{bmatrix}P_{11}\\ P_{21}\end{bmatrix}R,\qquad
P_{11}^TP_{11}+P_{21}^TP_{21}=I,
\]
with \(P_{11},P_{21}\) having at least as many rows as columns.  Show that
there exists an orthonormal \(Q\) such that
\[
A+\Delta A=QR,\qquad
\Delta A=F\Delta A_1+\Delta A_2,\qquad \lVert F\rVert_2\le 1 .
\]}
{Prove existence of an orthonormal factor and a combined data perturbation.}
{CS decomposition or a scoped substitute; polar/orthogonal factor facts.}
{Powerful but may require heavier exact linear algebra.}

\Exercise{Ch20 Ex. 20.2 -- Householder QR Least-Squares Backward Stability}
{Best}
{Least-squares backward stability}
{Prove the Householder QR least-squares stability theorem: if \(A\in
\mathbb R^{m\times n}\), \(m\ge n\), has full rank and
\(\min_x\lVert b-Ax\rVert_2\) is solved by Householder QR, then the computed
\(\widehat x\) is the exact least-squares solution to
\[
\min_x \lVert (b+\Delta b)-(A+\Delta A)x\rVert_2,
\]
with columnwise and right-hand-side bounds
\[
\lVert\Delta a_j\rVert_2\le \widetilde\gamma_{mn}\lVert a_j\rVert_2,\qquad
\lVert\Delta b\rVert_2\le \widetilde\gamma_{mn}\lVert b\rVert_2 .
\]}
{Show the computed LS solution exactly solves a perturbed LS problem with columnwise/RHS perturbation bounds.}
{Householder QR backward error; triangular solve backward error; LS normal-equation equivalence.}
{One of the best benchmark tasks.}

\Exercise{Ch20 Ex. 20.5 -- MGS Least-Squares Backward Stability}
{Best}
{MGS LS backward stability}
{For the modified Gram-Schmidt least-squares method applied as described in the
chapter, prove an analogue of the Householder QR theorem: the computed
\(\widehat x\) is the exact solution of a nearby least-squares problem, with
normwise or columnwise perturbation bounds of the same qualitative form as in
Theorem 20.3.}
{Formalize the MGS LS backward-stability result under a scoped MGS model.}
{MGS QR model; LS backward-error definitions.}
{Good if MGS algorithm support exists or is supplied in the task.}

\Exercise{Ch20 Ex. 20.6 -- Normal Equations Residual Error}
{Best}
{Residual error}
{For \(\min_x\lVert b-Ax\rVert_2\), let \(\widehat x\) be the solution computed by
the normal-equations method, let \(x\) be the exact solution, and define
\(\widehat r=b-A\widehat x\), \(r=b-Ax\).  Using the preceding normal-equations
error bounds, prove a first-order estimate of the form
\[
\lvert r^Tr-\widehat r^T\widehat r\rvert
\le c_{m,n}u\,\lVert x-\widehat x\rVert_2\lVert A\rVert_2
\bigl(\lVert A\rVert_2\lVert\widehat x\rVert_2+\lVert b\rVert_2\bigr)
+O(u^2).
\]}
{Derive a bound for \(\lvert r^T r-\widehat r^T\widehat r\rvert\).}
{Normal-equations backward error and residual algebra.}
{Strong because it is not just a direct theorem restatement.}

\Exercise{Ch20 Ex. 20.8 -- LS Backward Error with Fixed Right-Hand Side}
{Good}
{LS backward error}
{For the normwise least-squares backward error \(\eta_F(y)\) of Theorem 20.5,
evaluate the limiting case \(\theta=\infty\), i.e. the case \(\Delta b\equiv 0\)
where only \(A\) may be perturbed.}
{Derive the specialized expression or bound.}
{LS backward-error formula and singular-value/eigenvalue facts.}
{Good if the full formula has already been formalized.}

\Exercise{Ch20 Ex. 20.9 -- Stable LS Backward-Error Formula}
{Good}
{Stable formula}
{Prove the stable formula for the least-squares backward error:
\[
\eta_F(y)=\min\left\{\phi,\,
\sigma_{\min}\bigl([\,A\ \ \phi(I-rr^+)\,]\bigr)\right\},
\qquad
\phi=\sqrt{\mu}\,\frac{\lVert r\rVert_2}{\lVert y\rVert_2},
\]
where \(r=b-Ay\), \(\mu=\theta^2\lVert y\rVert_2^2/(1+\theta^2\lVert y\rVert_2^2)\),
and \(r^+\) denotes the pseudoinverse of \(r\) viewed as a vector.}
{Prove equivalence of the alternative formula to the original backward-error expression.}
{Singular-value formulas and projection residual algebra.}
{Good but may be exact-linear-algebra heavy.}

\Exercise{Ch20 Ex. 20.11 -- LSE First-Order Reduction}
{Good}
{First-order approximation}
{For the equality-constrained least-squares perturbation bound, specialize the
general formula to the unconstrained problem \(B=0\), \(d=0\).  Show that the
bound reduces, to first order, to the standard unconstrained least-squares
condition estimate.}
{Formalize the reduction under the unconstrained specialization.}
{First-order perturbation framework and LS definitions.}
{Useful if LSE definitions are available.}

\Exercise{Ch22 Ex. 22.8 -- Bidiagonal Inverse Perturbation}
{Best}
{Bidiagonal inverse perturbation}
{Let \(U\in\mathbb R^{n\times n}\) be upper bidiagonal and let
\(\Delta u_{ij}=\delta_{ij}u_{ij}\).  Derive the entrywise formula for
\((U+\Delta U)^{-1}-U^{-1}\) in terms of the diagonal and superdiagonal
\(\delta\)'s.  Then apply it to the perturbed bidiagonal factors \(U_k+\Delta U_k\)
from the Vandermonde analysis to prove a bound of the form
\[
(U_k+\Delta U_k)^{-1}=U_k^{-1}+F_k,\qquad
\lvert F_k\rvert\le \frac{(n+1)u}{1-(n+1)u}\lvert U_k^{-1}\rvert .
\]}
{Prove the exact formula and deduce the displayed bound.}
{Bidiagonal inverse formulas; product perturbation bounds.}
{Excellent structured-matrix stability task.}

\Exercise{Ch22 Ex. 22.11 -- Structured Vandermonde Condition Number}
{Best}
{Structured condition number}
{For the primal Vandermonde system \(Vx=b\), where
\(V=V(\alpha_0,\ldots,\alpha_n)\), define a structured condition number by
allowing relative perturbations \(\lvert\Delta\alpha_i\rvert\le
\epsilon\lvert\alpha_i\rvert\).  Show
\[
\operatorname{cond}(V)
=\frac{\lVert \lvert V^{-1}\operatorname{diag}(0,1,\ldots,n)V\rvert
\lvert x\rvert\rVert_\infty}{\lVert x\rVert_\infty},
\]
and derive the corresponding condition number for the dual system \(V^Ta=f\).}
{Formalize the first-order condition number formula for the primal system and optionally the dual analogue.}
{Vandermonde matrices; derivative of powers; componentwise perturbation.}
{Strong, but likely needs local Vandermonde definitions.}

\Exercise{Ch22 Ex. 22.3 -- Vandermonde Inverse Algorithm Stability}
{Good}
{Algorithm stability}
{For the Vandermonde inverse algorithm and its modified version, evaluate
left and right residuals of the computed inverses and compare them with GEPP.
As a scoped formal task, prove the deterministic sign/cancellation claim: the
algorithm performs subtractions of like-signed numbers, and use that fact to state
the resulting stability implication.}
{Use a scoped deterministic residual or subtraction-sign result, not the research subpart.}
{Vandermonde inverse algorithm definitions.}
{Avoid empirical or open-ended wording.}

\Exercise{Ch22 Ex. 22.7 -- Chebyshev-Vandermonde Conditioning}
{Good}
{Conditioning}
{For the Chebyshev-Vandermonde-like matrix
\(T=T(\alpha_0,\ldots,\alpha_n)\), prove the condition estimates
\[
\kappa_2(T)=\sqrt2
\quad\text{for}\quad
\alpha_i=\cos((i+\tfrac12)\pi/(n+1)),
\]
and
\[
\kappa_2(T)\le 2
\quad\text{for}\quad
\alpha_i=\cos(i\pi/n).
\]}
{Prove one of the two condition-number bounds.}
{Chebyshev polynomial identities and discrete orthogonality.}
{More exact approximation theory than FP stability, but condition-focused.}

\Exercise{Ch23 Ex. 23.6 -- 3M plus Strassen Error Bound}
{Best}
{Fast matrix multiply error}
{Combine the complex \(3M\) multiplication formula with Strassen multiplication
for the real products.  Prove the normwise error bound described after
Section 23.2.4: the computed complex product satisfies the Strassen bound
analogous to
\[
\lVert C-\widehat C\rVert
\le \left[\left(\frac n{n_0}\right)^{\log_2 12}(n_0^2+5n_0)-5n\right]
u\,\lVert A\rVert\lVert B\rVert+O(u^2),
\]
but with the extra constant multiplier and additive constant introduced by the
\(3M\) combination.}
{Derive the Strassen-style bound with the extra constants from the 3M combination.}
{Strassen normwise bound; 3M component formulas; complex matrix norms.}
{Excellent for showing normwise versus componentwise stability.}

\Exercise{Ch24 Ex. 24.1 -- FFT Convolution Integer Recovery}
{Best}
{FFT convolution error}
{For integer vectors \(x,y\), form the convolution coefficients \(z_k\) by padding
with zeros and using
\[
z=F_n^{-1}\bigl((F_nx)\cdot(F_ny)\bigr),
\]
where the dot denotes componentwise multiplication.  Derive a bound on
\(z-\widehat z\) for the FFT-based computation and give a sufficient condition
under which componentwise nearest-integer rounding of \(\widehat z\) recovers
the exact integer vector \(z\).}
{Prove that nearest-integer recovery is exact under a sufficient error bound.}
{DFT matrix norm/error model; componentwise multiplication error; inverse transform bound.}
{Strong but may require FFT/DFT local setup.}

\Exercise{Ch25 Ex. 25.1 -- Perturbed Contraction Iteration}
{Good}
{Perturbed iteration stability}
{Consider
\[
x_{k+1}=G(x_k)+e_k,\qquad \lVert e_k\rVert\le\alpha,
\]
where \(G:\mathbb R^n\to\mathbb R^n\) satisfies
\[
\lVert G(x)-a\rVert\le \beta\lVert x-a\rVert,\qquad 0<\beta<1.
\]
Show the ball of radius \(\alpha/(1-\beta)\) about \(a\) is invariant, that
outside this ball the distance to \(a\) strictly decreases, and that every
accumulation point lies in the ball.}
{Prove the invariant ball and accumulation-point bound.}
{Norm contraction algebra and recurrence bounds.}
{Good nonlinear stability task, less FP-specific.}

\Exercise{Ch27 Ex. 27.5 -- Robust Scaled 2-Norm}
{Good}
{Robust norm computation}
{Analyze the scaled BLAS \(x\)NRM2 algorithm for \(\lVert x\rVert_2\).  It keeps
\(t\) as the current scale and \(s\) as the scaled sum of squares, updating by
\[
\text{if }\lvert x_i\rvert>t:
\quad s:=1+s(t/x_i)^2,\ t:=\lvert x_i\rvert,
\qquad
\text{else } s:=s+(x_i/t)^2,
\]
and returns \(t\sqrt{s}\).  Prove the invariant establishing correctness and
state the robustness properties regarding overflow and underflow.}
{Formalize the invariant relating the scale and accumulated sum to \(\lVert x\rVert_2\).}
{Norm-square algebra; overflow/underflow-free exact model if included.}
{Good algorithmic task, but not central matrix stability.}

\Exercise{Ch27 Ex. 27.6 -- Pythagorean Sum Iteration}
{Good}
{Robust Pythagorean sum}
{For \(0<y_0\le x_0\), let \(p=\sqrt{x_0^2+y_0^2}\).  Show that Halley's method
applied to \(f(x)=x^2-p^2\) gives
\[
x_{n+1}=x_n\left(1+2\frac{p^2-x_n^2}{p^2+3x_n^2}\right),
\]
that \(x_n\le x_{n+1}\le p\), and, with \(y_n=\sqrt{p^2-x_n^2}\), derive the
scaled recurrence implemented by the pythag algorithm.  Also prove the cubic
error identity
\[
p-x_{n+1}=\frac{(p-x_n)^3}{p^2+3x_n^2}.
\]}
{Prove monotone/cubic convergence and the scaled recurrence.}
{Scalar iteration algebra and norm2/sqrt support.}
{Good robust-arithmetic task; keep the floating-point robustness claim scoped.}

\end{document}
